\chapter{Conclusion}\label{chap:conclusion}

To conclude, this memoire set out to explore how quantum computing challenges the foundations of modern cryptography, and why that matters. Each chapter tackled a different angle of this shift, helping me piece together the big picture: that adapting to disruptive technologies isn’t optional, it’s a must.

I started by diving into the fundamentals of quantum computing. Concepts like \gls{superposition} and \gls{entanglement} weren’t just theoretical curiosities, they explain how quantum computers can break things classical ones can’t. From there, I revisited classical cryptography, the kind that secures almost everything online today. Understanding how fragile that security becomes in a quantum context made the risk feel very real.

Looking into Shor’s and Grover’s algorithms showed just how deep the threat runs. Systems we trust daily, from banking logins to encrypted emails, could be rendered obsolete.

That’s why I found post-quantum cryptography (PQC) so important. These new cryptographic approaches are designed to survive quantum attacks, using hard mathematical problems that remain secure even in a quantum world. But exploring Chapter 7 made one thing very clear: switching to PQC is not easy. It’s technically complex, expensive, and forces organizations to rethink their infrastructure.

One of the key takeaways from this memoire is that resisting change just because it’s inconvenient is not an option when it comes to security. Yes, the transition to PQC comes with overhead, in time, money, and expertise. But the cost of doing nothing is far greater.

Personally, working on this topic has shown me just how interconnected technology and society are. The tools we build, or fail to update, have real-world consequences. This research has strengthened my belief that staying ahead in tech isn’t just about performance or innovation. Sometimes, it’s about resilience, preparation, and having the humility to change course when new realities emerge.

In short, this memoire succeeded in showing how critical it is to adapt, even when it’s hard. Quantum computing is coming, as I write advancements are being made, and the systems we rely on must evolve to meet it, not when it's convenient, but before it’s too late.
